\chapter{Inputs and outputs}\label{ch:c6}
This chapter is the reference for every command and every output, followed by a complete example set-up.

\begin{plain}
This is the look-up chapter: every command you can write in the input files, what it does and its default, and every output
file with its columns. There is no need to read it from start to finish; come here when you need the exact spelling of a
command or the meaning of an output column.
\end{plain}

\section{Input reference}\label{ch:inputs}
This chapter lists every command. Commands are case sensitive; separators may be commas or spaces, as in all Raven files.

\subsection{\texttt{.rvi} (main input file)}
\begin{table}[!htbp]\centering
\caption{Commands in the \cmd{.rvi} file.}
\begin{tabular}{P{0.34\linewidth}P{0.608\linewidth}}\toprule\hdr
Command & Meaning\\\midrule
\cmd{:GroundwaterModel MODFLOW6} & switches the coupling on (\cmd{NONE}: off, the default)\\
\cmd{:rvg\_Filename file} & name of the groundwater file, relative to the \cmd{.rvi} (default \cmd{<model>.rvg})\\
processes into \cmd{GROUNDWATER} & any process moving water into \cmd{GROUNDWATER} in a coupled HRU produces recharge; use
\cmd{:-->Conditional HRU\_GROUP} so uncoupled HRUs keep their own baseflow\\\bottomrule
\end{tabular}
\end{table}

\subsection{\texttt{.rvh} (HRUs and subbasins)}
The existing \cmd{AQUIFER\_PROFILE} column of \cmd{:HRUs} holds a profile name, or \cmd{[NONE]}. Reservoirs to be coupled
need \cmd{:SeepageParameters} and a lake HRU (\cmd{:HRUID}) with an \cmd{AQUIFER\_PROFILE}.

\subsection{\texttt{.rvp} (parameters)}
\begin{lstlisting}
:AquiferClasses
  :Attributes, K_HORIZ, K_VERT, SPEC_STORAGE, SPEC_YIELD, POROSITY
  :Units,      m/d,     m/d,    1/m,          none,       none
  GRAVEL,      50.0,    5.0,    1.0E-5,       0.25,       0.30
  CLAY_TILL,   0.001,   0.0001, 1.0E-4,       0.03,       0.40
:EndAquiferClasses
:AquiferProfiles
  # name, number of layers, then {class, thickness [m] or TO_BEDROCK, layer type} for each layer
  VALLEY_ALLUVIUM, 3, GRAVEL, 15.0, AQUIFER, CLAY_TILL, 5.0, AQUITARD, SAND, TO_BEDROCK, CONFINED_AQUIFER
  UPLAND_TILL,     1, UPLAND_TILL, TO_BEDROCK, AQUIFER
:EndAquiferProfiles
:AquiferProfileParameters
  :Attributes,     INITIAL_HEAD_DEPTH, SEEPAGE_LEAKANCE, RECHARGE_DELAY
  :Units,          m,                  1/d,              d
  VALLEY_ALLUVIUM, 3.0,                1.0,              0.0
:EndAquiferProfileParameters
\end{lstlisting}
\begin{table}[!htbp]\centering
\caption{Aquifer class attributes (\cmd{.rvp}).}
\begin{tabular}{P{0.27\linewidth}P{0.13\linewidth}P{0.528\linewidth}}\toprule\hdr
Aquifer class attribute & Default & Meaning\\\midrule
\cmd{K\_HORIZ} & required & horizontal hydraulic conductivity [m/d]\\
\cmd{K\_VERT} & $K_h/10$ & vertical hydraulic conductivity [m/d]\\
\cmd{SPEC\_STORAGE} & $10^{-5}$ & specific storage [1/m]\\
\cmd{SPEC\_YIELD} & 0.1 & specific yield [--]\\
\cmd{POROSITY} & 0.3 & total porosity [--] (reserved for transport)\\\bottomrule
\end{tabular}
\end{table}
Layer types: \cmd{AQUIFER} (or \cmd{UNCONFINED\_AQUIFER}), \cmd{CONFINED\_AQUIFER}, \cmd{AQUITARD} (or \cmd{CONFINING\_LAYER}),
\cmd{AQUICLUDE} (not as top layer). Classes must be defined before the profiles that use them.
\begin{table}[!htbp]\centering
\caption{Aquifer profile parameters (\cmd{.rvp}).}
\begin{tabular}{P{0.3\linewidth}P{0.13\linewidth}P{0.498\linewidth}}\toprule\hdr
Profile parameter & Default & Meaning\\\midrule
\cmd{INITIAL\_HEAD\_DEPTH} & 5 m & initial depth to water below land surface\\
\cmd{DEFAULT\_BEDROCK\_DEPTH} & 50 m & bedrock depth where no bedrock map is given\\
\cmd{SOIL\_ZONE\_DEPTH} & 0 m & model top below land surface\\
\cmd{SEEPAGE\_LEAKANCE} & 1 d$^{-1}$ & seepage conductance per m$^2$ of cell\\
\cmd{SEEPAGE\_SMOOTHING\_DEPTH} & 0.5 m & depth below the top where seepage starts\\
\cmd{RIVERBED\_K} & 0.5 m/d & riverbed hydraulic conductivity\\
\cmd{RIVERBED\_THICKNESS} & 1 m & riverbed thickness\\
\cmd{EXTINCTION\_DEPTH} & 0 m & water-table ET extinction depth (0 = off)\\
\cmd{RECHARGE\_DELAY} & 0 d & recharge delay time constant (0 = off)\\\bottomrule
\end{tabular}
\end{table}

\subsection{\texttt{.rvg} (groundwater file)}
\begin{longtable}{P{0.45\linewidth}P{0.498\linewidth}}
\caption{Commands in the \cmd{.rvg} file.}\\
\toprule\hdr
Command & Meaning (default)\\\midrule\endfirsthead
\caption[]{Commands in the \cmd{.rvg} file (continued).}\\
\toprule\hdr
Command & Meaning (default)\\\midrule\endhead
\cmd{:HRUGeometry file \{ID field\}} & HRU polygons, GeoJSON in lon/lat; ID field (\cmd{HRU\_ID}). \textbf{Required.}\\
\cmd{:RiverGeometry file \{subbasin field\}} & river lines; without a field, subbasins are assigned automatically\\
\cmd{:LandSurface file} & DEM, ESRI ASCII grid in lon/lat (HRU \cmd{ELEVATION})\\
\cmd{:BedrockSurface file} & bedrock elevation, same format (\cmd{DEFAULT\_BEDROCK\_DEPTH})\\
\cmd{:GridCellSize m} & cell size (1000 m)\\
\cmd{:MinCellCoverage f} & fraction of a cell that coupled HRUs must cover (0.25)\\
\cmd{:GridType REGULAR|QUADTREE|HRU\_MESH|FILE} & kind of grid (\cmd{REGULAR}); Section~\ref{sec:gridtypes}\\
\cmd{:MF6Simulation file} & couple an existing MODFLOW 6 simulation (Chapter~\ref{ch:existing})\\
\cmd{:MF6Model}, \cmd{:MF6StartDate} & model to couple (if several); date of the model's time zero\\
\cmd{:MF6CRS}, \cmd{:MF6GridFile}, \cmd{:OverlapWeights} & how HRUs are linked to the model's cells (one of them)\\
\cmd{:RavenTakesOver}, \cmd{:KeepPackages} & recharge/ET/UZF packages replaced by Raven or kept\\
\cmd{:StreamPackage p \{AUX a|DOMINANT\_HRU\}}, \cmd{:SeepageToRaven ON|OFF} & stream package to Raven's reaches; seepage package\\
\cmd{:GridRotation} $\theta$ & rotation of regular and quadtree grids, degrees counterclockwise (0)\\
\cmd{:GridRefinement RIVERS|WELLS|LAKES s}, \cmd{LINES|POLYGONS file s} & finer cells (size $s$ [m]) along rivers, around wells, in and around reservoir lakes, or along/around features in a file\\
\cmd{:GridFile file \{ID\}} & imported cell polygons, with \cmd{:GridType FILE} (ID field \cmd{CELL\_ID})\\
\cmd{:QuadtreeMaxLevel n} & maximum halvings of the base cell (from the refinement sizes)\\
\cmd{:FlowCorrection AUTO|XT3D|XT3D\_RHS|NONE} & XT3D conductance tensor (\cmd{AUTO}: off; \cmd{XT3D\_RHS} or \cmd{XT3D} to switch it on)\\
\cmd{:LayerConnection INDEX|ELEVATION} & connect layers by number or by overlapping elevation (\cmd{INDEX})\\
\cmd{:MF6Library path} & MODFLOW 6 library (else the environment variable \cmd{RAVEN\_MF6\_LIB}, next to the Raven executable, then \cmd{lib/mf6/}, then the system path; Section~\ref{sec:getmf6})\\
\cmd{:SeepageReturnTo SV} & Raven store receiving seepage (lowest soil layer; soil-less HRUs: \cmd{SURFACE\_WATER})\\
\cmd{:RiverBedDepth m | REFERENCE\_FLOW} & riverbed depth below the DEM (\cmd{REFERENCE\_FLOW})\\
\cmd{:GWInitialization STEADY\_STATE r} & steady-state start with recharge $r$ mm/d (off)\\
\cmd{:SurfaceWater\allowbreak Exchange SV group LEAKANCE L} & Raven store of an HRU group $\leftrightarrow$ aquifer\\
\cmd{:ReservoirExchange} & Raven reservoirs $\leftrightarrow$ aquifer\\
\cmd{:Wells} \dots\ \cmd{:EndWells} & table: ID, \cmd{NAME, LATITUDE, LONGITUDE, SCREEN\_TOP, SCREEN\_BOT}\\
\cmd{:ObservationWells} \dots & table: ID, \cmd{NAME, LATITUDE, LONGITUDE, SCREEN\_ELEV}\\
\cmd{:GeneralHeadBoundary name file HEAD h CONDUCTANCE c \{LAYERS ALL|TOP\}} & head-dependent boundary; $c$ per cell [m$^2$/d]\\
\cmd{:SpecifiedHeadBoundary name file HEAD h \{LAYERS ALL|TOP\}} & fixed-head boundary\\
\cmd{:SolverMax\allowbreak OuterIterations n} & Newton iterations per step (500)\\
\cmd{:SolverHead\allowbreak Tolerance m} & Newton head-change tolerance (0.00001 m)\\
\cmd{:HeadSaveFrequency n} & steps between saves of the full head field (30)\\
\cmd{:WriteNetCDFHeads} & write \cmd{GWHeads.nc} (Raven compiled with NetCDF)\\
\cmd{:LinkageCache file} & linkage cache location (\cmd{<output>/mf6/linkage.cache})\\\bottomrule
\end{longtable}
The block commands take one row per item, for example:
\begin{lstlisting}
:Wells
  :Attributes, NAME, LATITUDE, LONGITUDE, SCREEN_TOP, SCREEN_BOT
  1, PW_1, 59.27310, -124.73460, 1199.7, 1188.7
:EndWells
:ObservationWells
  :Attributes, NAME, LATITUDE, LONGITUDE, SCREEN_ELEV
  101, OW_NEAR, 59.27310, -124.73060, 1192.7
:EndObservationWells
:GeneralHeadBoundary   GHB1  boundary_line.geojson HEAD 1146.85 CONDUCTANCE 5000.0 LAYERS TOP
:SpecifiedHeadBoundary LAKE1 lake.geojson          HEAD 1226.5  LAYERS TOP
:SurfaceWaterExchange  DEPRESSION WetlandHRUs LEAKANCE 0.01
:ReservoirExchange
\end{lstlisting}

\subsection{\texttt{.rvt} (time series)}
\begin{lstlisting}
:WellRate 1 m3/d                  # well ID; negative = pumping
  1985-10-01 00:00:00 1.0 1826
  0.0
  ...
:EndWellRate
:BoundaryHead GHB1 m              # replaces HEAD of boundary GHB1
  1985-10-01 00:00:00 1.0 1826
  ...
:EndBoundaryHead
:ObservationData GW_HEAD 101 m    # observation-well ID; head AT each date
  1985-10-01 00:00:00 1.0 1826
  ...
:EndObservationData
\end{lstlisting}

\subsection{\texttt{.rvc} (initial conditions)}
\par\noindent\begin{minipage}{\linewidth}
\begin{lstlisting}
:InitialGWHeads DEPTH_BELOW_SURFACE 5.0     # or FROM_PROFILE (default) or ELEVATION z
\end{lstlisting}
\end{minipage}\par\smallskip
\cmd{:GWHeads nlay nrow ncol}, \cmd{:GWRechargeStore}, \cmd{:GWRiverLossCarry} and \cmd{:GWReservoirSeepage} are written by Raven into
\cmd{solution.rvc} for hotstart; they are read back automatically.

\section{Output reference}\label{ch:outputs}
All rows are stamped with the date at the \emph{end} of the step (the moment the values refer to).
\subsection{\texttt{GWBudget.csv} --- one row per step}
\begin{table}[!htbp]\centering
\caption{Columns of \cmd{GWBudget.csv}.}
\begin{tabular}{P{0.43\linewidth}P{0.518\linewidth}}\toprule\hdr
Column & Meaning\\\midrule
time, date & model time [d] and date at the end of the step\\
Raven recharge [m$^3$/d] & recharge leaving Raven (before any delay)\\
MF6 recharge [m$^3$/d] & recharge entering MODFLOW\\
seepage to Raven [m$^3$/d] & drain discharge at the model top\\
storage release [m$^3$/d] & water released from aquifer storage (negative: stored)\\
MF6 balance error [m$^3$/d], [\%] & $\varepsilon$ of Section~\ref{ch:balance}, and relative to gross flow\\
returned to Raven [m$^3$] & seepage volume added to Raven stores\\
outer iterations, converged & Newton iterations used; 1 target tolerance reached, 2 accepted at the relaxed tolerance, 0 not converged\\
recharge delay storage [m$^3$] & water in delay reservoirs\\
recharge onto CHD cells [m$^3$/d] & recharge of HRUs lying entirely on fixed-head cells\\
river leakage to aquifer, aquifer discharge to rivers [m$^3$/d] & river exchange, each direction\\
river exchange applied to Raven [m$^3$] & net volume actually added to (or taken from) Raven reaches\\
river loss carried to next step [m$^3$] & river loss not supplied in the step, taken from the river the next step (normally 0)\\
water-table ET [m$^3$/d] & evaporation from the water table\\
wells specified, wells actual [m$^3$/d] & requested and delivered well rates\\
GHB, CHD [m$^3$/d] & boundary flows (positive into the aquifer)\\
surface-store leakage / discharge [m$^3$/d], applied [m$^3$] & wetland and lake exchange\\
reservoir seepage sent, delivered [m$^3$/d], awaiting delivery [m$^3$] & reservoir exchange\\
reservoir gain not supplied by MODFLOW [m$^3$] & a gain the lake received that a drying cell could not supply; taken back from the water reaching the subbasin in the same step (normally 0)\\\bottomrule
\end{tabular}
\end{table}

\noindent One row of \cmd{GWBudget.csv} from the 20-year Liard run (selected columns):
\begin{center}\small
\begin{tabular}{ll}\toprule\hdr
Column & Value\\\midrule
\cmd{time [d]}, \cmd{date} & 3910, 1996-06-15\\
\cmd{Raven recharge [m3/d]} & 21\,365\,037.27\\
\cmd{MF6 recharge [m3/d]} & 21\,365\,037.27\\
\cmd{seepage to Raven [m3/d]} & 19\,087\,056.98\\
\cmd{storage release [m3/d]} & $-3\,925\,200.44$\\\bottomrule
\end{tabular}
\end{center}
Recharge leaving Raven equals recharge entering MODFLOW; on this wet June day the aquifer stores about 3.9 million m$^3$.

\subsection{Other files}
\begin{table}[!htbp]\centering
\caption{Other output files.}
\begin{tabular}{P{0.25\linewidth}P{0.698\linewidth}}\toprule\hdr
File & Content\\\midrule
\cmd{GWHRUState.csv} & per coupled HRU: overlap-weighted mean water table of its columns and mean depth to water\\
\cmd{GWHeads.csv} & head at each monitoring well, from $t=0$ (missing value $-1.2345$ when the screened cell is dry)\\
\cmd{GWConnectivity.csv} & number of separate wet areas; wet fraction of each profile\\
\cmd{mf6/grid\_cells.geojson} & every cell as a lon/lat polygon: number, active, profile, dominant HRU, area\\
\cmd{GWBudget.csv}, last column & with an existing model: net flow of the model's own packages (wells, boundaries, \dots) [m$^3$/d]\\
\cmd{GWModelSummary.txt} & MODFLOW 6 library and its version, grid type and size, MODFLOW format, HRU fidelity of the cells, active and pass-through cells, profiles, river/well/boundary/store/reservoir cells, linkage per HRU (areas, fraction linked, cells), whether the linkage came from the cache\\
\cmd{GWHeads.nc} & (optional) head(time, layer, y, x) with 2-D latitude and longitude, CF conventions\\
\cmd{mf6/} & the generated MODFLOW model, its list file (\cmd{gwf.lst}, \cmd{mfsim.lst}) and binary heads (\cmd{gwf.hds})\\
\cmd{Diagnostics.csv} & Raven's usual file, now also with \cmd{GW\_HEAD} statistics\\
\cmd{Watershed\allowbreak Storage.csv} & Raven's balance; includes all exchanged water\\
\cmd{solution.rvc} & includes \cmd{:GWHeads}, \cmd{:GWRechargeStore}, \cmd{:GWReservoirSeepage}\\\bottomrule
\end{tabular}
\end{table}

\section{A complete example set-up}
The following \cmd{.rvg} uses every option; lines not needed can be deleted.
\begin{lstlisting}
# ---- geometry ----
:HRUGeometry        my_hrus.geojson   HRU_ID
:RiverGeometry      my_rivers.geojson SubId
:LandSurface        dem_lonlat.asc
:BedrockSurface     bedrock_lonlat.asc
# ---- grid and solver ----
:GridCellSize       1000.0
:MinCellCoverage    0.25
:SolverHeadTolerance 0.00001
:SolverMaxOuterIterations 500
:HeadSaveFrequency  30
:WriteNetCDFHeads
:LinkageCache       cache/linkage.cache
# ---- coupling options ----
:MF6Library         /opt/mf6/libmf6.so
:SeepageReturnTo    SOIL[2]
:RiverBedDepth      REFERENCE_FLOW
:GWInitialization   STEADY_STATE 0.5
:SurfaceWaterExchange DEPRESSION WetlandHRUs LEAKANCE 0.01
:ReservoirExchange
# ---- stresses and observations ----
:Wells
  :Attributes, NAME, LATITUDE, LONGITUDE, SCREEN_TOP, SCREEN_BOT
  1, TownWell, 49.123, -80.456, 245.0, 230.0
:EndWells
:ObservationWells
  :Attributes, NAME, LATITUDE, LONGITUDE, SCREEN_ELEV
  101, OW_1, 49.120, -80.450, 240.0
:EndObservationWells
:GeneralHeadBoundary   East  east_edge.geojson HEAD 210.0 CONDUCTANCE 2000.0 LAYERS ALL
:SpecifiedHeadBoundary Lake  lake.geojson      HEAD 255.0 LAYERS TOP
\end{lstlisting}
with, in the \cmd{.rvt},
\begin{lstlisting}
:WellRate 1 m3/d
  2000-01-01 00:00:00 1.0 3
  -1500.0
  -1500.0
  -1800.0
:EndWellRate
:ObservationData GW_HEAD 101 m
  2000-01-01 00:00:00 1.0 3
  241.20
  241.18
  -1.2345
:EndObservationData
\end{lstlisting}
(\cmd{-1.2345} marks a missing observation, as elsewhere in Raven). Only three days are shown: like every Raven
forcing series, a \cmd{:WellRate} series must cover the whole simulation, and its time interval must divide evenly
into one day (Raven stops with a message otherwise). Observation series may be shorter.
